% paper/main.tex
% AAAI 2027 Submission
% Target: Main Track — Machine Learning

\documentclass[letterpaper]{article}
\usepackage{times}
\usepackage{helvet}
\usepackage{courier}
\usepackage{graphicx}
\usepackage{amsmath}
\usepackage{amssymb}
\usepackage{amsthm}
\usepackage{booktabs}
\usepackage{xcolor}
\usepackage{hyperref}
\usepackage{algorithm}
\usepackage{algpseudocode}
\usepackage[margin=1in]{geometry}

% Theorem environments
\newtheorem{definition}{Definition}
\newtheorem{hypothesis}{Hypothesis}

% Custom commands
\newcommand{\rki}{\textsc{RKI}}
\newcommand{\bacc}{\textsc{BACC}}
\newcommand{\forget}{\mathcal{D}_f}
\newcommand{\retain}{\mathcal{D}_r}
\newcommand{\TODO}[1]{\textcolor{red}{\textbf{[TODO: #1]}}}

\title{Snowballing at the Forgetting Boundary: \\
How Partial Unlearning Destabilizes Chain-of-Thought Reasoning}

\author{
  Anonymous Author(s) \\
  Affiliation \\
  \texttt{email@domain.com}
}

\begin{document}
\maketitle

%% -----------------------------------------------
\begin{abstract}
\TODO{Write after experiments confirm hypothesis.}
Machine unlearning has emerged as a critical technique for removing knowledge from
large language models (LLMs). However, we identify an overlooked failure mode:
partial unlearning destabilizes chain-of-thought (CoT) reasoning. We show that
residual knowledge traces near the forget-set boundary act as corrupted nodes
in the reasoning chain, causing errors to compound---a phenomenon we term
\textit{forgetting-boundary snowballing}. We formalize this via the
Residual Knowledge Interference (\rki) score, a mutual-information measure
of forget-concept residue as a function of CoT depth. Mechanistic analysis
reveals that \TODO{N} attention heads in layers \TODO{L1--L2} become unstable
post-unlearning and propagate errors downstream. To mitigate this, we propose
\bacc{} (Boundary-Aware CoT Checkpointing), a lightweight inference-time
intervention requiring no retraining. Experiments across three models and three
unlearning methods demonstrate \TODO{X\%} reduction in snowballing with
negligible utility loss.
\end{abstract}

%% -----------------------------------------------
\section{Introduction}
\label{sec:intro}

Machine unlearning~\cite{cao2015towards} has gained significant attention as a
means of selectively erasing knowledge from trained models in response to privacy
regulations and safety concerns. Recent work has established effective methods
for gradient-based forgetting~\cite{chundawat2023bad,zhang2024npo}, yet a critical
downstream effect has gone unstudied: \textit{what happens to multi-step reasoning
after unlearning?}

\TODO{Fill: paragraph on CoT snowballing setup}

\TODO{Fill: paragraph stating hypothesis}

We make the following contributions:
\begin{enumerate}
    \item We identify and characterize \textit{forgetting-boundary snowballing},
    a novel failure mode where partial unlearning amplifies CoT errors.
    \item We introduce the \rki{} score, an information-theoretic metric
    quantifying residual knowledge interference in reasoning chains.
    \item We provide mechanistic evidence of which layers and attention heads
    drive this instability via probing analysis.
    \item We propose \bacc{}, an inference-time mitigation requiring no
    additional model training.
\end{enumerate}

%% -----------------------------------------------
\section{Related Work}
\label{sec:related}

\subsection{Machine Unlearning}
Early work on machine unlearning focused on classification
models~\cite{thudi2022unrolling}. Chundawat et al.~\cite{chundawat2023bad}
proposed the Bad Teacher method, using an incompetent teacher to induce
targeted forgetting. \TODO{Add NPO, TOFU, recent surveys.}

\subsection{Chain-of-Thought Reasoning}
Wei et al.~\cite{wei2022chain} introduced CoT prompting as a method for
eliciting step-by-step reasoning. Snowballing in CoT---where errors compound
across reasoning steps---was formalized by Gan et al.~\cite{gan2025snowball}
using an information-theoretic framework. \TODO{Add faithfulness, mechanistic interp refs.}

\subsection{Mechanistic Interpretability}
\TODO{Add Elhage circuits, probing classifier refs.}

%% -----------------------------------------------
\section{Problem Formulation}
\label{sec:formulation}

\subsection{Setup}
Let $\mathcal{M}$ be a language model trained on dataset $\mathcal{D} = \forget \cup \retain$,
where $\forget$ is the forget set and $\retain$ is the retain set.
After applying an unlearning algorithm $\mathcal{U}$, we obtain $\mathcal{M}_u = \mathcal{U}(\mathcal{M}, \forget)$.

\subsection{CoT Reasoning Chains}
A reasoning chain of depth $D$ is a sequence of steps
$\mathbf{c} = (c_1, c_2, \ldots, c_D)$ where each step $c_d$ is generated
autoregressively conditioned on prior steps. Let $h_l^{(d)}$ denote the
hidden state at layer $l$ during generation of step $c_d$.

\subsection{Residual Knowledge Interference (\rki) Score}
\begin{definition}[\rki{} Score]
The Residual Knowledge Interference score at layer $l$ and chain depth $d$ is:
\begin{equation}
    \rki(l, d) = I\!\left(h_l^{(d)};\, \mathcal{C}_{\forget}\right)
    \label{eq:rki}
\end{equation}
where $\mathcal{C}_{\forget}$ is the forget-concept representation
(centroid of $\forget$ embeddings) and $I(\cdot;\cdot)$ denotes mutual information.
\end{definition}

\begin{hypothesis}
When $\rki(l, d) > \tau$ for some threshold $\tau$ and chain depth $d > k$,
downstream CoT steps exhibit compounding error accumulation.
\end{hypothesis}

\TODO{Add formal error cascade definition.}

%% -----------------------------------------------
\section{Experiments}
\label{sec:experiments}

\subsection{Setup}
\textbf{Models.} We evaluate on Qwen2.5-1.5B-Instruct~(1.5B parameters),
loaded in 4-bit NF4 quantization. All unlearning is performed via LoRA adapters
($r{=}8$, $\alpha{=}16$, targeting \texttt{q\_proj} and \texttt{v\_proj}).

\textbf{Unlearning Methods.} We compare Gradient Ascent (GA)~\cite{thudi2022unrolling},
Bad Teacher~\cite{chundawat2023bad}, and NPO~\cite{zhang2024npo}.

\textbf{Dataset.} We use the TOFU benchmark~\cite{maini2024tofu}, which contains
fictitious author biographies. The forget set (\texttt{forget01}) contains 40 QA pairs
about authors to be unlearned; the retain set (\texttt{retain99}) contains 3{,}960 QA pairs.

\textbf{Unlearning Verification.} GA unlearning with 5 epochs yields a forget-set
perplexity increase from 12.73 to 1{,}146.32 (90$\times$), confirming effective
behavioral forgetting at the output level.

\subsection{RQ1: Does Snowballing Increase Post-Unlearning?}
\TODO{Fill with snowball measurement results (measure\_snowball.py).}

\subsection{RQ2: Mechanistic Analysis --- Where Does Residual Knowledge Live?}

We train linear probes at 8 evenly-spaced layers to classify whether a hidden-state
activation originated from the forget set or the retain set. Table~\ref{tab:rki}
reports probe accuracy for the base model and the GA-unlearned model.

\begin{table}[h]
\centering
\caption{Linear probe accuracy (forget vs.\ retain classification) per layer,
before and after Gradient Ascent unlearning on Qwen2.5-1.5B.}
\label{tab:rki}
\begin{tabular}{lcc}
\toprule
Layer & Base Model & Unlearned Model \\
\midrule
0  (embedding)  & 0.583 & \textbf{1.000} \\
4               & 1.000 & 1.000 \\
8               & 1.000 & 1.000 \\
12              & 1.000 & 1.000 \\
16              & 1.000 & 1.000 \\
20              & 1.000 & 1.000 \\
24              & 1.000 & 1.000 \\
28 (final)      & \textbf{1.000} & 0.958 \\
\bottomrule
\end{tabular}
\end{table}

\textbf{Key finding:} Unlearning primarily affects the final layer (probe accuracy
drops from 1.000 to 0.958), while all intermediate layers retain perfect separability.
Notably, the embedding layer (layer~0) becomes \emph{more} separable after unlearning
(0.583$\to$1.000), suggesting that gradient ascent reorganizes early representations
while only suppressing output-level behavior. This confirms that GA unlearning is
\textit{surface-level}: internal representations of the forget concept remain intact,
creating the conditions for residual knowledge interference during multi-step reasoning.

\subsection{RQ3: Does \bacc{} Mitigate Snowballing?}
\TODO{Fill with Phase 4 results.}

%% -----------------------------------------------
\section{The \bacc{} Method}
\label{sec:bacc}

\TODO{Fill after Phase 4 implementation.}

Algorithm sketch:
\begin{algorithm}
\caption{Boundary-Aware CoT Checkpointing (\bacc{})}
\begin{algorithmic}[1]
\Require Unlearned model $\mathcal{M}_u$, probe $P$, steering vector $\mathbf{v}$, threshold $\tau$
\For{each CoT step $d = 1, \ldots, D$}
    \State Generate hidden state $h^{(d)}$ at monitored layer $l^*$
    \If{$P(h^{(d)}) > \tau$}  \Comment{Boundary detected}
        \State Apply steering: $h^{(d)} \leftarrow h^{(d)} + \alpha \mathbf{v}$
    \EndIf
    \State Generate token $c_d$ from modified $h^{(d)}$
\EndFor
\end{algorithmic}
\end{algorithm}

%% -----------------------------------------------
\section{Conclusion}
\label{sec:conclusion}

\TODO{Fill last.}

%% -----------------------------------------------
\bibliographystyle{plain}
\bibliography{references}

\end{document}
